\section{Proofs for the complete graph}\label{app:complete}

This appendix proves Theorems~\ref{thm:simplex-crossing} and~\ref{thm:simplex-balance} and
Corollary~\ref{cor:simplex-modes}. We use the notation of Section~\ref{sec:completegraph}, so
$Q=J-dI$, $\sigma=p-h$, $u=h/p$ and $v=h/\sigma$. The first lemma writes the law in a second
form, with blocks between refreshes in place of runs.

\begin{lemma}[the refresh form]\label{lem:refresh}
Let $N\ge1$ and let $n$ be a composition with support $F$. Sums over $m$ run over the integer vectors
with $1\le m_i\le n_i$ for $i\in F$, and $M=\sum_im_i$.
\begin{enumerate}[(a)]
\item With $0^0=1$,
\[
P_N(n)=\sum_m\Bigl(\sum_{i\in F}\mu_im_i\Bigr)\frac{(M-1)!}{\prod_im_i!}\,
h^{M-1}\sigma^{N-M}\prod_{i\in F}\binom{n_i-1}{m_i-1}.
\]
\item Let $1/d<p<1$, so that $0<u<1$ and $v>0$. Then
\begin{equation}\label{eq:simplex-positive-sum}
S_n(u)=(1-u)^{N-1}\sum_mM!\,v^{M-1}\prod_{i\in F}\frac1{m_i!}\binom{n_i-1}{m_i-1},
\end{equation}
which is a sum of positive terms, and
\begin{equation}\label{eq:simplex-balance-integral}
S_n(u)=\frac{(1-u)^{N-1}}{v}\int_0^\infty e^{-t}\prod_{i\in F}B_{n_i}(vt)\,dt,\qquad
B_n(y)=\sum_{m=1}^n\binom{n-1}{m-1}\frac{y^m}{m!}.
\end{equation}
\end{enumerate}
\end{lemma}

\begin{proof}
(a) The transition matrix is $I+hQ=\sigma I+hJ$, so a word $x_1\cdots x_N$ has probability
$\mu_{x_1}\prod_{s<N}(\sigma\mathbf 1\{x_s=x_{s+1}\}+h)$. We expand each factor into its two
terms and call the steps where the term $h$ is chosen the refreshes. If there are $M-1$
refreshes, they cut the word into $M$ blocks, and a term of the expansion is not zero only if the
word is constant on each block. So $P_N(n)$ is the sum, over the sequences of $M$ block labels
and over the block lengths, of $\mu_{\ell_1}h^{M-1}\sigma^{N-M}$, where $\ell_1$ is the label of
the first block. If $m_i$ blocks carry the label $i$, there are $m_i(M-1)!/\prod_jm_j!$ sequences
of labels that start with $i$. Consecutive labels may be equal. Given the labels, the lengths of
the $m_i$ blocks of the label $i$ form an ordered sum of $n_i$ with $m_i$ positive parts, in
$\binom{n_i-1}{m_i-1}$ ways. Multiplying and summing gives (a).

(b) For the uniform start $\sum_i\mu_im_i=M/d$. By~\eqref{eq:SnVN} we divide (a) by
$V_N=p^{N-1}/d$, and $h^{M-1}\sigma^{N-M}/p^{N-1}=u^{M-1}(1-u)^{N-M}=(1-u)^{N-1}v^{M-1}$, since
$\sigma/p=1-u$. This gives~\eqref{eq:simplex-positive-sum}. We write
$M!=\int_0^\infty e^{-t}t^M\,dt$ and sum over each $m_i$ separately. All terms are nonnegative, so
the sum and the integral can be exchanged, and this gives~\eqref{eq:simplex-balance-integral}.
\end{proof}

When $\sigma\ge0$, part (a) has a direct reading. The chain keeps its state with probability
$\sigma$ and otherwise draws a fresh state uniformly from all $d$ states, possibly the same one.
The proof does not need $\sigma\ge0$. At $p=1/d$ we have $\sigma=0$ and only $M=N$ survives,
which gives the multinomial law for the uniform start.

\subsection{Proof of Theorem~\ref{thm:simplex-crossing}}

Let $F$ be a face of $k$ states and $n$ a balanced composition of $F$, so that $P_N(n)/V_N=S_n(u)$
by~\eqref{eq:SnVN}, with $u=(1-p)/((d-1)p)$.

\emph{Part (i), one crossing.} To count $W_n(j)$, choose a sequence of $j+1$ run
colors in which successive colors differ. If color $i$ appears $m_i$ times, its
run lengths can be chosen in $\binom{n_i-1}{m_i-1}$ ways. So every coefficient
of $S_n$ is nonnegative, and the first nonzero one is $W_n(k-1)=k!$, from the
$k!$ orders of $k$ single runs. Hence $S_n$ increases strictly from $0$ to
$\infty$ on $u>0$ and equals $1$ exactly once. As $p$ increases over $(0,1)$,
$u$ decreases over $(0,\infty)$. So $P_N(n)=V_N$ at exactly one $p_{N,k}\in(0,1)$, and the composition is
more likely than a vertex below $p_{N,k}$ and less likely above it.

\emph{Part (i), monotonicity in $N$.} Choose the balanced compositions for increasing
$N$ by adding $1$ to one of the smallest coordinates. Each coefficient $W_n(j)$
is then nondecreasing. To see this, map a word to the word with an extra copy of
the incremented color just before its first occurrence. This keeps the number
of switches, and it is injective, since deleting the first letter of the first
run of that color inverts it. The coefficient of $u^k$ is
\[
W_n(k)=\frac{k!(k-1)}2\sum_i(n_i-1)=\frac{k!(k-1)}2(N-k).
\]
To see this, note that exactly one run color appears twice. For each choice of it there are
$(k+1)!/2-k!=k!(k-1)/2$ sequences of run colors in which successive colors
differ (all orders, minus those with the two copies adjacent), and $n_i-1$ ways to
split its $n_i$ visits into two runs. This coefficient increases strictly with
$N$. So at every $u>0$ the value $S_n(u)$ increases strictly with $N$, the root
of $S_n(u)=1$ decreases strictly with $N$, and $p=1/(1+(d-1)u)$ increases.

\emph{Part (i), the first crossing.} If $N=k$, every coordinate is $1$ and
$S_n(u)=k!u^{k-1}$, which gives $p_{k,k}$. Since the root decreases with $N$,
every root has $u\le(k!)^{-1/(k-1)}<1$, and $u<1$ means $h<p$, that is,
$p>1/d$.

\emph{Part (i), the face maximum.} For $p\ge1/d$, Theorem~\ref{thm:simplex-balance} gives
$M_F=P_N(n)$. For $p<1/d$ we have $u>1$, which exceeds the root of $S_n(u)=1$, so
$M_F\ge P_N(n)>V_N$. Hence $M_F=V_N$ exactly at $p_{N,k}$, with the same two inequalities.

\emph{Part (ii).} By Proposition~\ref{prop:nextterm}(ii) with $K_F=D_k/d$, $\lambda_F=d-k$ and
$c_F=-(k-1)/8$ (Corollary~\ref{cor:paperB}), for large $N$ and uniformly for $1\le T\le N^{1/3}$,
\[
\log M_F=-(k-1)L+\tfrac{k-1}2\log T+\log(D_k/d)-(d-k)T-\frac{k-1}{8T}+O(T^{-2}+T^2/N).
\]
By
\eqref{eq:vertexmax}, $\log V_N=-\log d-(d-1)T+O(T^2/N)$. Subtracting, and using
$\log D_k=-(k-1)a_k$, gives the display. With $T=L-\frac12\log L+t$ and $t$ in a compact set,
$T-L+\frac12\log T=t+\frac12\log(T/L)=t+O(\log L/L)$, and the terms $1/(8T)$ and the error tend
to $0$.

\emph{Part (iii).} Let $G$ be a vertex, so that $M_G=V_N$. By
Theorem~\ref{thm:crossing}(i) the function $D_N=\log(M_F/M_G)$ has a zero in $[L/2,2L]$, since
$\Delta=\delta=k-1$, and by (i) this zero is $T_{N,k}$. By Corollary~\ref{cor:paperB},
$\log(K_F/K_G)=\log D_k=-(k-1)a_k$ and $c_F-c_G=-(k-1)/8$. So~\eqref{eq:crossingnext} reads
$T(k-1)(L-\frac12\log T-T+a_k)=-\frac{k-1}8+O(1/L)$ at $T=T_{N,k}$, and dividing by
$(k-1)T$ gives the implicit equation. For the explicit form put
$z_*=L-\frac12\log L+a_k+e_L/L$ with $e_L=\frac14\log L-\frac12a_k+\frac18$.
Since $\log z_*=\log L-\frac{\log L}{2L}+\frac{a_k}L+O(\frac{(\log L)^2}{L^2})$,
we get $z_*+\frac12\log z_*-1/(8z_*)=L+a_k+O((\log L)^2/L^2)$. The derivative of
$z+\frac12\log z-1/(8z)$ is at least $1$ for $z>0$, so
$T_{N,k}-z_*=O((\log L)^2/L^2)$, which is~\eqref{eq:simplex-root-expansion}.

\emph{Part (iv).} We have $a_k=\frac12\log(4\pi)-f(k)$ with
$f(x)=(x+\frac12)\log x/(x-1)$. For $x\ge2$ the sign of $f'(x)$ is the sign of
$\Delta(x)=2x^2-x-1-3x\log x$. Here $\Delta(2)=5-6\log2>0$, and
$\Delta'(x)=4x-4-3\log x>0$ for $x\ge2$, since $\Delta'(2)>0$ and
$\Delta''(x)=4-3/x>0$. So $f$ increases strictly and $a_k$ decreases strictly. This order is
checked in Lean. \qed

\subsection{Proof of Theorem~\ref{thm:simplex-balance}}

Let $\mu_a=\mu_0$ for $a\in F$. By~\eqref{eq:runform}, $P_N(n)=\mu_0p^{N-1}S_n(u)$ for every composition
with support $F$, so it is enough to show that the move strictly increases $S_n(u)$. The
polynomial $S_n$ is symmetric in the coordinates of $n$, so all balanced compositions of $F$ have the same
value, and repeated moves reach a balanced composition from any composition with support $F$. So the move gives
the theorem. At $p=1/d$ we have $u=1$ and $S_n(1)=N!/\prod_in_i!$, the number of words with the
counts $n$, and the move multiplies it by $n_b/(n_a+1)>1$. Let $1/d<p<1$, so that $\sigma>0$ and
$v>0$.

\emph{Log-concavity.} We show that for each $y>0$ the sequence $B_n(y)$,
$n\ge1$, is strictly log-concave. Write
\[
\frac{B_n(y)}y=b_{n-1},\qquad
b_j=\sum_{i=0}^j\binom ji a_i,\quad
a_i=\frac{y^i}{(i+1)!},\qquad
c_j=\sum_{i=0}^j\binom ji a_{i+1}.
\]
Pascal's identity gives $b_{j+1}=b_j+c_j$. The ratio $c_j/b_j$ is the mean of the
strictly decreasing sequence $f_i=a_{i+1}/a_i=y/(i+2)$ under the weights
$\lambda_i=\binom jia_i$, $0\le i\le j$, normalized. Let $\lambda'_i=\binom{j+1}ia_i$,
$0\le i\le j+1$, be the weights for $j+1$, and put $\lambda_{j+1}=0$. On
$0\le i\le j$ the ratio $\lambda'_i/\lambda_i=(j+1)/(j+1-i)$ increases in $i$. So for
$i<l$ we have $\lambda'_i\lambda_l-\lambda'_l\lambda_i\le0$, with strict inequality when
$l=j+1$. Hence
\[
\sum_i\lambda'_if_i\sum_l\lambda_l-\sum_i\lambda'_i\sum_l\lambda_lf_l
=\sum_{i<l}(\lambda'_i\lambda_l-\lambda'_l\lambda_i)(f_i-f_l)<0,
\]
since $f_i>f_l$ for $i<l$. So the mean decreases strictly, $c_{j+1}/b_{j+1}<c_j/b_j$,
and $b_{j+1}/b_j=1+c_j/b_j$ decreases strictly with $j$. This is the strict
log-concavity. Since $B_n(y)=(y/n)L^{(1)}_{n-1}(-y)$, it is a Tur\'an-type
inequality for the Laguerre polynomials $L^{(1)}_{n-1}(-y)$.

\emph{The transfer.} Let $n_b\ge n_a+2$. By log-concavity,
$B_{n_a+1}(y)/B_{n_a}(y)>B_{n_b}(y)/B_{n_b-1}(y)$, that is,
$B_{n_a+1}(y)B_{n_b-1}(y)>B_{n_a}(y)B_{n_b}(y)$. We apply this
for $y=vt>0$ in~\eqref{eq:simplex-balance-integral} and integrate. So the move
strictly increases $S_n(u)$. \qed

\subsection{Proof of Corollary~\ref{cor:simplex-modes}}

For large $N$, $T\le CL$ forces $h\to0$ and so $p>1/d$. By
Theorem~\ref{thm:simplex-balance} every global mode is a vertex or a balanced
composition, and all balanced compositions with the same number of positive coordinates have the
same probability. So we compare the ratios $M_F/V_N$, $2\le\lvert F\rvert\le d$, which equal $S_n(u)$ at
the balanced compositions $n$, with $1$ and with each other. They depend on $F$ only through $k=\lvert F\rvert$.
Fix $0<c<1$.

\emph{Case $T\le cL$.} The ratio $S_n(u)$ increases with $u$, and $u$ increases
with $T$. So it is at most its value at $T=cL$, which is
$\exp(-(k-1)(1-c)L+O(\log L))\to0$ by Theorem~\ref{thm:simplex-crossing}(ii). So the
vertices are the only modes, and $p>p_{N,d}$.

\emph{Case $cL\le T\le CL$.} Theorem~\ref{thm:simplex-crossing}(ii) gives
$\log(M_F/V_N)=(k-1)(y-a_k)+o(1)$ for $\lvert F\rvert=k$, with $y=T-L+\frac12\log T$, uniformly. The number $y$
is the same for every $k$. Fix $\delta>0$, with
$\delta<(a_{d-1}-a_d)/2$ when $d\ge3$. This is possible since $a_2>\cdots>a_d$.
If $y\le a_d-\delta$, every face ratio is at most
$e^{-(k-1)\delta+o(1)}<1$, so only the vertices are modes. If
$y\ge a_d+\delta$, the full center beats the vertices, and it beats a
balanced composition with $2\le k<d$ positive coordinates by
\[
(d-k)(y-a_d)+(k-1)(a_k-a_d)+o(1)>0 .
\]
If $|y-a_d|<\delta$, a proper face with $k\ge2$ states has
$(k-1)(y-a_k)+o(1)\le-(k-1)\delta+o(1)<0$, since $a_k-a_d>2\delta$. So it
loses to the vertices, and the comparison of the vertices with the full centers
is decided by the unique crossing $p_{N,d}$ of Theorem~\ref{thm:simplex-crossing}(i).
In each subcase the conclusion agrees with (i) to (iii).

The last claim follows from~\eqref{eq:simplex-root-expansion} and
$p_{N,k}=1-(d-1)T_{N,k}/N$, since $a_2>\cdots>a_d$. \qed
